\mwnewpart
\renewcommand\runningtitle{Section 2. Installing RetrOS/LT}
\chapter{Installing RetrOS/LT}
\sectioncontents

\section{Two Ways to Use RetrOS/LT}

\begin{itemize}
\item \textbf{From diskettes.} Start the computer from the DOS diskette, or
  from a copy of it, and run your programs from the diskettes. Nothing
  needs installing. To make a startup diskette of your own, see ``Making a
  Disk Start the Computer'' in Section~4.
\item \textbf{From the fixed disk.} Install puts RetrOS/LT on the first
  fixed disk, drive C, so that the computer starts from it. This Section
  describes Install.
\end{itemize}

\section{Before You Install}\ixs{Install}\ixs{INSTALL command}

\begin{itemize}
\item \textbf{Copy the two diskettes}, and install from the copies.
\item \textbf{Save what is on the fixed disk.} If the disk holds files you
  want, copy them to diskettes first. Install can keep the files on a disk
  that is already formatted, but partitioning or formatting a disk erases
  it.
\item \textbf{A new fixed disk} must have been low-level formatted, as its
  maker describes, before Install can use it.
\end{itemize}

Install asks before it changes the fixed disk. You can leave it at any
question by pressing \key{Esc}; nothing further is changed.

\section{Starting Install}

\begin{enumerate}
\item Put the DOS diskette in drive A and close the drive door.
\item Turn the computer on, or, if it is on, press
  \key{Ctrl}\,+\,\key{Alt}\,+\,\key{Del}.
\item Enter the date and the time when RetrOS/LT asks for them.
\end{enumerate}

Install then starts:

\begin{dialogue}
\ \ RetrOS/LT Version 1.0 Install

\ \ Welcome to RetrOS/LT.

\ \ Install puts RetrOS/LT on the fixed disk of this computer,
\ \ drive C:, so that the computer starts from it.

\ \ Before any change is made to the fixed disk you are asked first.
\ \ You can leave Install at any question by pressing Esc.

\ \ Press Enter to continue, or Esc to leave Install.
\end{dialogue}

Press \key{Enter}. Install then asks you to accept the RetrOS/LT License,
which is printed in Appendix~D and is on the diskette as
\file{LICENSE.TXT}:

\begin{dialogue}
\ \ Do you accept the terms of the RetrOS/LT License (Y/N)?
\end{dialogue}

Press \key{Y} to go on. If you press \key{N}, Install ends and nothing is
changed. To read the license first, press \key{N} and type
\typed{TYPE LICENSE.TXT | MORE}. To start Install again, type
\typed{INSTALL}.

If the computer has no fixed disk, Install says so and ends. RetrOS/LT
runs from the diskettes as they are.

\section{A Disk with No DOS Partition}\ixs{partition}\ixs{FDISK command}

A fixed disk must be divided into partitions before it can be used. If
yours has no DOS partition, Install says so:

\begin{dialogue}
\ \ The fixed disk has no DOS partition.

\ \ A fixed disk must be partitioned before it can be used.

\ \ \ \ W\ \ \ Use the whole disk for RetrOS/LT.
\ \ \ \ \ \ \ \ Anything on the disk now will be lost.

\ \ \ \ F\ \ \ Run FDISK, to set up the partitions yourself.

\ \ \ \ Esc Leave Install without changing the disk.

\ \ Press W, F or Esc.
\end{dialogue}

\begin{itemize}
\item Press \key{W} to give the whole disk to one DOS partition, as
  \typed{FDISK} does when you answer \typed{Y} to ``Do you wish to use the
  entire fixed disk for DOS''. This is the usual choice.
\item Press \key{F} to run \typed{FDISK} yourself, if you want to leave
  part of the disk for another operating system. See ``Partitioning a
  Fixed Disk'' in Section~4.
\end{itemize}

A DOS 2 partition holds at most 65,535 sectors, about 32 MB. On a larger
disk the rest is left unused.

The computer must then restart. Leave the diskette in drive A and press
any key. RetrOS/LT starts again, asks the date and time, and Install goes
on from where it stopped. Accept the license once more.

\section{Choosing the Directory}\ixs{directory, install}

Install asks where to put the commands and device drivers:

\begin{dialogue}
\ \ Where should RetrOS/LT be installed?

\ \ The programs and drivers go in one directory of drive C:.
\ \ C:\bs{}DOS is the usual place. To use another, type over the name
\ \ below. The files that start the computer always go in C:\bs.

\ \ C:\bs{}DOS

\ \ Press Enter to accept, or Esc to leave Install.
\end{dialogue}

Press \key{Enter} to use \file{C:\bs{}DOS}. To use another directory, press
\key{Backspace} to erase the name, type the new one, and press
\key{Enter}. Each name in the path may have up to eight letters or digits;
a path such as \file{C:\bs{}SYS\bs{}DOS} is allowed, and Install makes each
directory in it.

\section{Formatting Drive C}\ixs{formatting drive C}

If drive C is not formatted, Install says that it will format it and
waits for \key{Enter}. It then runs \typed{FORMAT C:/S}, which asks you to
press a key to begin. Formatting checks every sector of the disk and takes
some minutes.

If drive C is already formatted, Install asks what to do:

\begin{dialogue}
\ \ Drive C: is already formatted and may hold files.

\ \ \ \ K\ \ \ Keep the files on drive C: and add RetrOS/LT to them.
\ \ \ \ \ \ \ \ An existing CONFIG.SYS and AUTOEXEC.BAT are kept as
\ \ \ \ \ \ \ \ CONFIG.OLD and AUTOEXEC.OLD.

\ \ \ \ F\ \ \ Format drive C:. Everything on it will be lost.

\ \ \ \ Esc Leave Install without changing the disk.

\ \ Press K, F or Esc.
\end{dialogue}

\begin{itemize}
\item Press \key{K} to keep your files. Install runs \typed{SYS C:}, which
  puts the RetrOS/LT system files in place of those of the system that was
  there.
\item Press \key{F} to format the disk. Install asks once more,
  \shown{Format drive C: (Y/N)?}, before it does.
\end{itemize}

\caution{\typed{SYS} can put the system files on a disk only if the first
two entries of its root directory are free or already hold system files,
and the start of the disk has room for them. A disk formatted by DOS 2 or
by RetrOS/LT has that room; a disk formatted by a later version of DOS may
not. If \typed{SYS} cannot do it, Install says so and stops. Copy your
files to diskettes, run Install again, and choose to format drive C.}

\section{Copying the Files}

Install copies the files of the DOS diskette to the directory you chose,
showing each name, and puts \file{COMMAND.COM} in \file{C:\bs} as well.
It then asks for the second diskette:

\begin{dialogue}
\ \ Insert the RetrOS/LT Supplemental Programs diskette in
\ \ drive A:.

\ \ Press Enter to continue, or Esc to leave Install.
\end{dialogue}

Change the diskettes and press \key{Enter}. Install checks that it is the
right diskette before it copies anything from it.

\section{The Startup Files}\ixs{CONFIG.SYS}\ixs{AUTOEXEC.BAT}

Last, Install writes \file{C:\bs{}CONFIG.SYS} and
\file{C:\bs{}AUTOEXEC.BAT}. If they were there already, the old ones are
kept as \file{CONFIG.OLD} and \file{AUTOEXEC.OLD}. With the usual
directory, they are:

\begin{dialogue}
FILES=20
BUFFERS=15
DEVICE=C:\bs{}DOS\bs{}ANSI.SYS
\end{dialogue}

\begin{dialogue}
DATE
TIME
PROMPT \$P\$G
PATH C:\bs{}DOS
\end{dialogue}

On a computer with less than 384 KB, \typed{BUFFERS=10} is written
instead, to leave more memory for programs. The Personal Computer and the
XT keep no time while they are off, so \file{AUTOEXEC.BAT} asks the date
and time each time the computer starts. The prompt shows the current
directory, and the commands can be used from any directory. Section~5
describes these files.

When Install is done, it says so:

\begin{dialogue}
\ \ RetrOS/LT is installed on drive C:.

\ \ Remove the diskette from drive A:, then press Enter to restart
\ \ the computer from the fixed disk.
\end{dialogue}

Open the drive door and press \key{Enter}. The computer starts from the
fixed disk.

\section{If Install Stops}

If something goes wrong, Install says what it was, and that nothing
further was changed. The usual causes:

\begin{itemize}
\item \shown{The fixed disk could not be read}: the adapter or its cable
  is not seated, or the disk has not been low-level formatted.
\item \shown{Drive C: could not be formatted}: run \typed{FORMAT C:/S}
  from the DOS diskette to see the message.
\item \shown{A file could not be copied to drive C:}: the disk is full, or
  the diskette is damaged. Try another copy of the diskette.
\end{itemize}

Install works out where it stopped from the state of the disk, so you can
simply run it again.
